% исправления 27.05.2005
\chapter{Теории и модели}

\section{Аксиомы равенства}
        \label{equality-axioms}

Пусть сигнатура~$\sigma$ включает в себя двуместный предикат
равенства (записываемый традиционно $x=y$). Интерпретация этой
сигнатуры называется нормальной,%
\index{Нормальная интерпретация}%
\index{Интерпретация!нормальная}
если предикат равенства
интерпретируется как тождественное совпадение элементов
носителя.

Возникает естественный вопрос. Пусть имеется некоторая
теория~$T$ (множество замкнутых формул) в языке, сигнатура
которого включает равенство. Мы знаем что теория имеет модель
(интерпретацию, в которой все формулы из~$T$ истинны) тогда и
только тогда, когда она непротиворечива. В каком случае она
имеет нормальную модель (нормальную интерпретацию, в которой все
формулы из~$T$ истинны)?

Чтобы ответить на этот вопрос, введём аксиомы равенства. Пусть
$\sigma$\т произвольная сигнатура. \emph{Аксиомами равенства}%
\index{Аксиомы!равенства|defin}%
\index{Равенства аксиомы|defin}
в сигнатуре~$\sigma$ будут формулы
        \begin{align*}
\forall x&\,(x=x),\\
\forall x\forall y &\,((x=y)\to(y=x)),\\
\forall x\forall y\forall z&\, (((x=y)\land (y=z))\to(x=z))
        \end{align*}
(называемые аксиомами рефлексивности, симметричности и
        \index{Рефлексивность}%
        \index{Симметричность}%
        \index{Транзитивность}%
транзитивности). Это ещё не всё. Для каждого функционального
символа мы формулируем аксиому равенства, которая говорит, что
его значение не меняется, если аргументы заменить на равные.
Например, для двуместного функционального символа~$f$
соответствующая аксиома выглядит так:
       $$
\forall x_1 \forall x_2 \forall y_1 \forall y_2\,
 (((x_1=x_2)\land(y_1=y_2))
\to (f(x_1,y_1)=f(x_2,y_2))).
       $$
Для предикатных символов аксиомы равенства говорят, что
истинный предикат остаётся истинным, если заменить аргументы на
равные. Например, для двуместного предикатного символа~$A$ аксиома
такова:
       $$
\forall x_1 \forall x_2 \forall y_1 \forall y_2\,
 (((x_1=x_2)\land(y_1=y_2)\land A(x_1,y_1)
                    \to A(x_2,y_2)).
       $$
(Нет необходимости специально говорить, что предикат, соответствующий
символу~$A$, остаётся
ложным при замене аргументов на равные, так как равенство
симметрично.)

\begin{theorem}[полноты для нормальных моделей]
\index{Теорема!о полноте для нормальных моделей|defin}%
        \label{normal-models}%
Тео\-рия~$T$ сигнатуры~$\sigma$ с равенством имеет нормальную
        \index{Нормальная модель}%
модель тогда и только тогда, когда она остаётся непротиворечивой
после добавления аксиом равенства.
        %
\end{theorem}

\begin{proof}
        %
Прежде всего заметим, что теоремы о корректности и полноте
(раздел~\ref{predicate-completeness})
позволяют говорить о совместности вместо непротиворечивости.

В нормальной модели теории~$T$ аксиомы равенства
истинны, так что в одну сторону утверждение теоремы очевидно.
Нам осталось показать, что если теория $T$ совместна с аксиомами
равенства, то она имеет нормальную модель.

Возьмём произвольную интерпретацию, в которой истинны формулы
из~$T$ и аксиомы равенства. Пусть $M$\т её носитель. В этой
интерпретации предикат $=$ не обязан быть настоящим равенством;
он представляет собой некоторое бинарное отношение на~$M$.
Поскольку выполнены аксиомы равенства, это отношение
рефлексивно, симметрично и транзитивно (является отношением
эквивалентности). Следовательно, множество~$M$ разбивается на
классы эквивалентности; множество этих классов обозначим $M'$
(его можно назвать фактор\д множеством $M$ по данному отношению
эквивалентности). Класс элемента~$x$ будем обозначать $[x]$.

Аксиомы равенства позволяют корректно определить интерпретацию c
носителем~$M'$. В самом деле, истинность аксиомы для
функционального символа~$f$ (приведённой выше в качестве
примера) гарантирует, что класс $[f(x,y)]$ зависит лишь от
классов $[x]$ и~$[y]$, но не от выбора~$x$ и~$y$ внутри класса.
Аналогичным образом аксиомы для предикатных символов позволяют
корректно определить предикаты на классах.

Полученная интерпретация с носителем~$M'$ по построению
нормальна. Осталось убедиться, что в ней истинны те же самые
формулы, что и в~$M$ (в том числе все формулы теории~$T$). Это
почти очевидно с интуитивной точки зрения:~$M$ отличается
от~$M'$ лишь тем, что каждый элемент представлен несколькими
равноправными копиями, которые со всех точек зрения ведут себя
одинаково.

Формально говоря, мы доказываем, что формула~$\varphi$ истинна в
интерпретации $M$ на оценке $\pi$ тогда и только тогда, когда
$\varphi$ истинна в $M'$ на оценке $\pi'$, при которой значение
любой переменной $\xi$ есть класс, содержащий значение
переменной $\xi$ при оценке~$\pi$. Это легко сделать индукцией
по построению формулы $\varphi$.
        %
\end{proof}

\begin{problem}
        %
Покажите, что из аксиом равенства для сигнатуры $\sigma$
выводится формула
        $$
 \varphi \land (x=y) \to \varphi(y/x),
        $$
если подстановка в правой части корректна. (Указание:
это очевидно следует из теоремы о полноте, но можно провести
и чисто синтаксическое рассуждение индукцией по построению
формулы~$\varphi$.)
        %
\end{problem}

\begin{problem}
        \label{clones}%
        \index{Увеличение мощности}%
Покажите, что если теория~$T$ (не обязательно с равенством)
имеет модель мощности~$\alpha$, то она имеет и модель
любой большей мощности. (Указание: элементы модели можно
\лк клонировать\пк\ в любом количестве.)
        %
\end{problem}

Из теоремы о полноте для нормальных моделей легко следует аналог
теоремы о компактности
(теорема~\ref{predicate-compactness-theorem},
с.~\pageref{predicate-compactness-theorem}) для нормальных
моделей.

\begin{theorem}[компактности для нормальных моделей]
        \label{normal-compactness}%
        \index{Теорема!компактности для нормальных моделей}%
        \index{Компактность}%
Если всякое конечное подмножество теории $T$ в сигнатуре
с равенством имеет нормальную модель, то и теория~$T$
имеет нормальную модель.
        %
\end{theorem}

\begin{proof}
        %
Любое конечное подмножество теории $T$ остаётся непротиворечивым
при добавлении аксиом равенства (поскольку имеет нормальную модель).
Значит, и вся теория $T$ остаётся непротиворечивой при
добавлении аксиом равенства (вывод противоречия использует
конечное число формул) и потому имеет нормальную модель.
        %
\end{proof}

\begin{problem}
        \index{Линейный порядок}
        \index{Частичный порядок, продолжение до линейного}
Применив теорему о компактности, докажите, что всякий частичный
порядок может быть продолжен до линейного. (Указание. Рассмотрим
частично упорядоченное множество как модель теории, в сигнатуре
которой есть равенство, порядок и константы для всех элементов
множества, а формулами являются равенства и неравенства между
константами. Добавим к ней утверждение о сравнимости любых двух
элементов. Покажите, что любое конечное множество формул
полученной теории непротиворечиво, используя тот факт, что
частичный порядок на конечном множестве продолжается до
линейного.)
        %
\end{problem}

\begin{problem}
        \index{Присоединение корней многочлена}%
        \index{Алгебраическое замыкание}%
        %
Используя теорему о компактности, докажите, что для всякого поля
$k$ сушествует его расширение $k'$, в котором всякий многочлен с
коэффициентами из~$k$ имеет корень. (Указание. Утверждение о
существовании корня у многочлена с данными коэффициентами можно
записать в виде формулы. Любое конечное множество таких формул
совместно с аксиомами поля, так как можно по очереди
присоединить корни соответствующих многочленов.)
        %
\end{problem}

\begin{problem}
        %
Пусть $\Gamma$\т множество замкнутых формул в сигнатуре с
равенством. Покажите, что замкнутая формула $\varphi$ этой
сигнатуры истинна во всех нормальных моделях $\Gamma$ тогда и
только тогда, когда она выводима из $\Gamma$ и аксиом равенства.
        %
\end{problem}

Утверждение последней задачи является аналогом
теоремы~\ref{semantics-provability}
(с.~\pageref{semantics-provability}) для теорий с равенством.
        \index{Теория!с равенством}%
        \index{Аксиомы!равенства}%
        \index{Равенство}%
Иногда вообще рассматривают только такие теории. При этом
равенство является обязательным элементом сигнатуры, аксиомы
равенства (их число зависит от сигнатуры) считаются частью
исчисления предикатов, а интерпретации рассматриваются только
нормальные. При этом теория имеет [нормальную] модель тогда и
        \index{Нормальная модель}%
только тогда, когда она непротиворечива [вместе с аксиомами
равенства]; формула выводима из теории~$\Gamma$ [и аксиом
равенства] тогда и только тогда, когда она верна во всех
[нормальных] моделях теории~$\Gamma$ \итп (в квадратных скобках
указаны подразумеваемые слова).

\section{Повышение мощности}
        \label{lowenheim-skolem-2}
        \index{Повышение мощности}%
Теорема Левенгейма\ч Сколема
      \glossary{\Лёвенгейм}%
      \glossary{\Сколем}%
(теорема~\ref{lowenheim-skolem-theorem} в
разделе~\ref{lowenheim-skolem-1}) позволяла уменьшать
мощность интерпретации (она утверждала, что для любой
бесконечной интерпретации конечной или счётной сигнатуры
существует элементарно эквивалентная ей счётная подструктура). В
этом разделе мы рассмотрим обратную задачу\т расширение
интерпретации до элементарно эквивалентной интерпретации большей
мощности. Соответствующее утверждение также называют теоремой
Левенгейма\ч Сколема.

Прежде всего отметим, что без требования нормальности такое
утверждение бессодержательно: как уже говорилось на
с.~\pageref{clones}, мы можем дублировать элементы интерпретации в
произвольном количестве. Поэтому мы предполагаем, что все
рассматриваемые интерпретации нормальны (равенство
интерпретируется как тождественное совпадение).

\begin{theorem}[Левенгейма\ч Сколема о повышении мощности]
\index{Теорема!Левенгейма\ч Сколема о повышении мощности|defin}%
        \label{lowenheim-skolem-theorem-2}%
Пусть $A$\т бесконечная нормальная интерпретация некоторой
сигнатуры $\sigma$ с равенством. Тогда существует нормальная
интерпретация $B\supset A$ сколь угодно большой мощности,
        \index{Элементарное расширение}%
являющаяся элементарным расширением~$A$.
        %
\end{theorem}

(Это означает, согласно определению на
с.~\pageref{elementary-extension}, напомним, что интерпретация
предикатных и функциональных символов в~$B$ продолжает их
интерпретацию в~$A$ и что формулы сигнатуры $\sigma$, параметрам
которых приданы значения из~$A$, одновременно истинны в~$A$ и
в~$B$.)

\begin{proof}
        %
Сформулируем утверждение теоремы в терминах теорий и моделей.
Пусть $A$\т произвольная нормальная интерпретация
сигнатуры~$\sigma$. Рассмотрим сигнатуру~$\sigma\!_A$, которая
получается из~$\sigma$ добавлением констант\т по одной для
каждого элемента множества~$A$. Эта сигнатура имеет естественную
нормальную интерпретацию с носителем~$A$: значением каждой
константы является соответствующий ей элемент. (Возможно, что в
$\sigma$ изначально было достаточно констант и всякий
элемент~$A$ был значением некоторой константы. Тогда эта
процедура лишняя, но и вреда от неё нет.)

Рассмотрим теорию $\Th_A(A)$, состоящую из формул
сигнатуры~$\sigma\!_A$, истинных в $A$ при указанной
интерпретации. Всякое элементарное расширение~$B$
интерпретации~$A$ будет моделью теории $\Th_A(A)$. В самом деле,
зам\-к\-ну\-тая формула $\varphi(a_1,\dots,a_n)$ сигнатуры~$\sigma\!_A$
получается подстановкой констант $a_1,\dots,a_n$ вместо
параметров в какую\д то формулу
$\varphi(x_1,\dots,x_n)$ сигнатуры~$\sigma$. (Мы используем не
вполне корректные обозначения, в частности, отождествляем
элементы $a_1,\dots,a_n$ множества~$A$ с константами для них.)
Её истинность в $B$ (или в~$A$) равносильна истинности
формулы~$\varphi(x_1,\dots,x_n)$ при значениях параметров
$x_1\hm\mapsto a_1,\dots,x_n\hm\mapsto a_n$\т фор\-маль\-но говоря,
следует воспользоваться леммой~2 на~с.~\pageref{trouble-lemma}.
Поэтому по определению элементарного расширения все формулы
из~$\Th_A(A)$ будут истинны и в $B$.

Верно и обратное: любая нормальная
модель теории $\Th_A(A)$ естественно
определяет элементарное расширение интерпретации~$A$. В самом
деле, пусть дана нормальная модель этой теории с носителем~$B$.
Тогда каждый элемент множества
$A$ (точнее, соответствующая этому элементу константа)
интерпретируется
некоторым элементом множества~$B$. Разным элементам множества~$A$
соответствуют разные элементы в~$B$, так как формула $a_1\ne
a_2$, истинная в $A$, должна быть истинной и в~$B$. Таким
образом, $A$ вкладывается в $B$ и можно отождествить его с
некоторым подмножеством множества~$B$. Это отождествление
корректно в том смысле, что предикаты и функциональные символы
интерпретируются согласованным образом. В самом деле, атомарные
формулы вида $P(a_1,\dots,a_n)$, а также формулы $\lnot
P(a_1,\dots,a_n)$ и $f(a_1,\dots,a_n)\hm= a$, истинные в $A$,
истинны и в~$B$. Истинные в~$A$ формулы вида
$\varphi(a_1,\dots,a_n)$ принадлежат $\Th_A(A)$ и потому истинны
и в~$B$; ложные в~$A$ формулы имеют отрицания в $\Th_A(A)$ и
потому ложны в~$B$.

Таким образом, для доказательства теоремы Левенгейма\ч Сколема о
повышении мощности осталось построить нормальную модель теории
$\Th_A(A)$, имеющую сколь угодно большую мощность. Это можно
сделать так: добавим множество новых констант~$c_i$ и формулы
$c_i\ne c_j$ (для всех $i\hm\ne j$) к теории $\Th_A(A)$.
Полученная теория будет совместной по теореме компактности для
нормальных моделей. (В самом деле, любая конечная часть её имеет
нормальную модель, поскольку содержит конечное число новых
констант, и им можно придать различные значения в~$A$.) Поэтому
и вся теория имеет нормальную модель. Всем константам $c_i$
соответствуют в этой модели разные элементы (поскольку истинна
формула~$c_i\hm\ne c_j$), поэтому мощность этой модели может
быть сколь угодно большой, если использовать достаточно много
констант.

Этот же приём будет использован нами при построении
нестандартной модели арифметики
(с.~\pageref{nonstandard-n-construction}).
        %
\end{proof}

Приведённое рассуждение даёт оценку мощности снизу. Можно
получить и в точности нужную мощность:

\begin{theorem}
        \label{infinite-spectrum}%
Пусть $A$\т бесконечная нормальная интерпретация
сигнатуры~$\sigma$ (с равенством) и пусть $\beta$\т мощность, не
меньшая мощностей сигнатуры~$\sigma$ и интерпретации~$A$. Тогда
существует нормальное элементарное расширение~$B\hm\supset A$
мощности $\beta$.
        %
\end{theorem}

\begin{proof}
        %
Мощность сигнатуры~$\sigma_A$ есть максимум из мощностей
$\sigma$ и $A$; после добавления новых констант в
количестве~$\beta$ штук получится сигнатура мощности $\beta$, и
согласно теореме~\ref{model-cardinality}
(с.~\pageref{model-cardinality}) найдётся модель
множества~$\Th_A(A)$ мощности $\beta$. Преобразование её в
нормальную модель (факторизация) может лишь уменьшить мощность,
но $\beta$ различных элементов у нас заведомо есть.
        %
\end{proof}

Аналогичный приём (добавление констант) позволяет легко доказать
такое утверждение:

\begin{theorem}
        %
        \index{Повышение мощности}%
Если теория (в произвольной сигнатуре с равенством) имеет сколь
угодно большие конечные нормальные модели, то она имеет и
бесконечную нормальную модель.
        %
\end{theorem}

\begin{proof}
        %
Добавим к теории бесконечное число новых констант и аксиомы о
том, что все они различны. Любой конечный фрагмент расширенной
теории имеет нормальную модель (возьмём достаточно большую
конечную модель и проинтерпретируем в ней константы). По теореме
компактности и вся расширенная теория имеет нормальную модель,
которая и будет бесконечной нормальной моделью исходной теории.
        %
\end{proof}

Вообще можно задать себе такой естественный вопрос. Пусть есть
некоторая теория (или даже просто одна формула). Каковы могут
быть мощности её нормальных моделей? Как мы видели, для теорий с
конечной сигнатурой верно одно из двух: либо бесконечных моделей
вовсе нет, либо есть бесконечные модели всех мощностей. Это
гарантируют теоремы Левенгейма\ч Сколема об элементарной
подмодели (теорема~\ref{lowenheim-skolem-theorem}) и о повышении
мощности (теорема~\ref{lowenheim-skolem-theorem-2}).

Что можно сказать про мощности конечных моделей? Для каждой
формулы рассмотрим множество всех возможных мощностей её
конечных моделей. Его иногда называют \emph{спектром}%
\index{Спектр формулы|defin}%
\index{Формула!спектр|defin}
формулы.
Это множество может быть устроено довольно сложным образом:
например, для формулы, выражающей аксиомы поля, спектр состоит
из всех степеней простых чисел.

\begin{problem}
        (\textbf{а})~%
Укажите формулу, спектр которой состоит из всех чётных положительных
чисел.
        (\textbf{б})~%
Укажите формулу, спектр которой состоит из всех нечётных чисел.
        (\textbf{в})~%
Укажите формулу, спектр которой состоит из всех составных чисел.
        %
\end{problem}

Любопытно, что \emph{проблема конечного спектра}%
\index{Проблема конечного спектра|defin}
(стоящая в
книге Кейслера и Чэна~\cite{keisler-chang} под номером~1 среди \лк старых проблем
теории моделей\пк), неожиданно оказалась связана с центральной
проблемой теории сложности вычислений\т так называемой \лк
проблемой перебора\пк. (Проблема конечного спектра состоит в
следующем: верно ли, что дополнение (до~$\mathbb{N}$) к спектру
любой формулы является спектром некоторой другой формулы?)

В качестве примера использования теоремы о повышении мощности
докажем теорему Гильберта о нулях
(теорема~\ref{nullstellensatz}), не проводя элиминацию
кванторов. Пусть система уравнений имеет решение в поле~$k'$,
являющемся расширением алгебраически замкнутого поля~$k$.
Покажем, что она имеет решение и в~$k$. Построим элементарное
расширение $k''\hm\supset k$ очень большой мощности. Теперь $k'$
можно вложить в $k''$ (это вложение строится по трансфинитной
рекурсии: добавляя алгебраический элемент, мы пользуемся
алгебраической замкнутостью~$k''$, добавляя трансцендентный
элемент, мы пользуемся большой мощностью~$k''$). Значит, система
имеет решение в $k''$. Поскольку $k''$ было элементарным
расширением, то система имеет решение в~$k$.

Другое любопытное применение теоремы о повышении мощности
таково. Назовём линейно упорядоченное множество $M$
\emph{однородным},%
\index{Однородное множество|defin}%
\index{Множество!однородное|defin}
если для любых двух возрастающих
последовательностей $x_1\hm<x_2\hm<\ldots\hm<x_n$ и
$y_1\hm<y_2\hm<\ldots\hm<y_n$ найдётся автоморфизм
множества~$M$, переводящий $x_i$ в~$y_i$ (при всех
$i\hm=1,\ldots,n$).

\begin{theorem}
        %
Для всякой бесконечной мощности найдётся однородное линейно
упорядоченное множество такой мощности.
        %
\end{theorem}

\begin{proof}
        %
Множество рациональных чисел (и вообще любое счётное плотное
линейно упорядоченное множество) однородно. В самом деле,
соответствие между двумя наборами его элементов постепенно
продолжается до автоморфизма (добавляем
элементы поочерёдно с той или другой стороны). Другой способ
убедиться в этом\т вспомнить о том, что это рациональные числа,
и взять кусочно\д линейный автоморфизм.

Для каждого $n$ фиксируем способ продолжения автоморфизмов с
$n$\д элементных подмножеств в виде функции $f_n$ с $2n+1$
аргументами: $f(z,x_1,\dots,x_n,y_1,\dots,y_n)$ означает
элемент, в который переходит $z$ при автоморфизме, переводящем
$x_i$ в $y_i$. Рассмотрим теперь $\mathbb{Q}$ как интерпретацию
сигнатуры, включающей порядок и все $f_i$. По теореме о
повышении мощности можно найти элементарно эквивалентную
интерпретацию любой заданной мощности. Поскольку свойства
функций $f_n$ выражаются формулами, получится однородное
линейно упорядоченное множество заданной мощности.
        %
\end{proof}

\section{Полные теории}

В этом разделе мы попытаемся систематизировать уже известные нам
понятия и факты.

\begin{itemize}
        \item
Начнём с напоминаний. \emph{Сигнатурой}%
\index{Сигнатура|defin}
мы называли набор
предикатных и функциональных символов. Среди формул данной
сигнатуры выделяют \emph{замкнутые}%
\index{Замкнутая формула|defin}%
\index{Формула!замкнутая|defin}
(формулы без параметров).
Сигнатура имеет \emph{интерпретации},%
\index{Интерпретация|defin}%
\index{Сигнатура!интерпретация|defin}
в которых замкнутые
формулы этой сигнатуры бывают \emph{истинными}%
\index{Истинная формула|defin}%
\index{Формула!истинная|defin}
и \emph{ложными}.%
\index{Ложная формула|defin}%
\index{Формула!ложная|defin}
Произвольное множество замкнутых формул данной сигнатуры
называется \emph{теорией}%
\index{Теория|defin}
в этой сигнатуре. \emph{Моделью}%
\index{Модель!теории|defin}%
\index{Теория!модель|defin}
теории называется интерпретация, в которой все формулы теории
истинны. Теория называется \emph{совместной},%
\index{Совместная теория|defin}%
\index{Теория!совместная|defin}
если она имеет модель.

        \item
Теория называется \emph{теорией с равенством}, если она включает
        \index{Теория!с равенством}%
в себя аксиомы равенства (а её сигнатура содержит символ
равенства). Интерпретация теории с равенством называется
\emph{нормальной}, если равенство интерпретируется как
        \index{Нормальная интерпретация}%
        \index{Интерпретация!нормальная}%
совпадение элементов носителя интерпретации. Совместная теория с
равенством имеет нормальную модель (получаемую из произвольной
        \index{Модель!нормальная}%
модели факторизацией по отношению равенства).

        \item
Говорят, что замкнутая формула $\varphi$ \emph{выводима в теории
$T$} (является \emph{теоремой теории~$T$}),
     \index{Выводимая формула|defin}%
     \index{Формула!выводимая в теории|defin}%
     \index{Теорема!теории|defin}%
     \index{Теория!теорема|defin}%
если формула $\varphi$ получается из аксиом исчисления предикатов
и формул теории $T$ по правилам вывода. (Обозначение: $T\vdash\varphi$.)

Формула $\varphi$ выводима в теории~$T$ тогда и только тогда,
когда в исчислении предикатов выводится некоторая формула вида
$\tau\to\varphi$, где $\tau$\т конъюнкция конечного числа формул
из~$T$.

Формула~$\varphi$
\emph{семантически следует} из~$T$, если она истинна в любой
        \index{Семантическое следование}%
        \index{Следование семантическое}%
модели теории~$T$ (обозначение: $T\hm\vDash\varphi$). Семантическое
следование равносильно выводимости
(теорема~\ref{semantics-provability},
с.~\pageref{semantics-provability}). Взяв в качестве $\varphi$
тождественно ложную формулу $\perp$ (скажем, отрицание
тавтологии), приходим к понятиям \emph{противоречивости}%
\index{Противоречивость теории|defin}
($T\hm\vdash\perp$) и \emph{несовместности}%
\index{Несовместность теории|defin}
(${T\vDash\perp}$,
$T$ не имеет моделей).
В противоречивой теории
выводима любая формула (соответствующей сигнатуры).

        \item
Непротиворечивая теория~$T$ \emph{полна}%
\index{Полная теория|defin}%
\index{Теория!полная|defin}
(в данной сигнатуре), если для
любой замкнутой формулы~$\varphi$ этой сигнатуры либо формула
$\varphi$, либо её отрицание $\lnot\varphi$ выводится из~$T$.

        \item
Для произвольной интерпретации~$M$ произвольной
сигнатуры~$\sigma$ можно рассмотреть \emph{элементарную теорию
интерпретации}~$M$,%
\index{Элементарная теория интерпретации|defin}
обозначаемую $\Th(M)$ и состоящую из всех
истинных в~$M$ замкнутых формул сигнатуры~$\sigma$. Очевидно,
эта теория полна (одна из формул $\varphi$ и $\lnot\varphi$ ей
принадлежит). Две интерпретации $M_1$ и $M_2$ \emph{элементарно
эквивалентны},%
\index{Интерпретации!элементарно эквивалентные|defin}
если их элементарные теории совпадают: $\Th(M_1)=\Th(M_2)$.
        \item
Теория~$T$ называется \emph{конечно аксиоматизируемой},%
\index{Конечно аксиоматизируемая теория|defin}%
\index{Теория!конечно аксиоматизируемая|defin}
если существует
конечное множество $T'$ теорем теории~$T$, из которых выводятся
все утверждения из~$T$ (другими словами, если существует
конечная теория, имеющая то же самое множество теорем).
        \item
Теория с равенством, имеющая конечную или счётную сигнатуру,
называется \emph{категоричной в счётной мощности},%
      \index{Теория!категоричная в счётной мощности|defin}
если все её счётные нормальные модели изоморфны.
Категоричность в данной несчётной мощности определяется
аналогично.

        \item
Теория с конечной сигнатурой называется~\emph{разрешимой},%
\index{Разрешимая теория|defin}%
\index{Теория!разрешимая|defin}
если существует алгоритм, который по произвольной замкнутой
формуле определяет, выводима ли она в этой теории или нет.
        %
\end{itemize}

\begin{problem}
        %
Покажите, что добавление к теории любой её теоремы не меняет
множества теорем.
        %
\end{problem}

Прежде чем переходить к примерам, сделаем два простых наблюдения.

\begin{theorem}[критерий Лося\ч Воота]
\index{Критерий!Лося\ч Воота|defin}%
\glossary{\Лось}%
\glossary{\Воот}%
        \label{categorical-complete}%
Не\-про\-ти\-во\-ре\-чи\-вая теория~$T$ с равенством
в конечной или счётной
сигнатуре, не имеющая конечных моделей и
категоричная в счётной мощности, полна.
        %
\end{theorem}

\begin{proof}
        %
Предположим, что ни одна из формул $\varphi$ и $\lnot\varphi$ не
выводима в теории $T$. Тогда обе теории $T\cup\{\lnot\varphi\}$
и $T\cup\{\varphi\}$ непротиворечивы. По
теореме~\ref{countable-model} (с.~\pageref{countable-model}) они
имеют счётные модели, которые остаются счётными после
факторизации (перехода к нормальным моделям), поскольку
теория~$T$ не имеет конечных моделей. Эти счётные модели должны
быть изоморфными (в силу категоричности). С~другой стороны, в
одной из них истинна формула $\lnot\varphi$, а в другой\т
формула~$\varphi$, так что они даже не элементарно эквивалентны
(мы знаем из раздела~\ref{elementary-equivalence}, что такого
быть не может).
        %
\end{proof}

Аналогично доказывается и общая форма критерия Лося\ч Воота:

\begin{theorem}
\index{Критерий!Лося\ч Воота!общая форма|defin}%
        \label{general-vaught}%
Непротиворечивая теория с равенством в конечной или счётной
сигнатуре, не имеющая конечных моделей и категоричная в данной
несчётной мощности~$\alpha$, полна.
        %
\end{theorem}

\begin{proof}
        %
Пусть теория~$T$ не полна и к ней можно присоединить без
противоречия любую из формул~$\varphi$ и $\lnot\varphi$.
Рассмотрим счётные нормальные модели теорий
$T\hm\cup\{\varphi\}$ и $T\hm\cup\{\lnot\varphi\}$. По
теореме~\ref{infinite-spectrum} увеличим их мощности до~$\alpha$
и получим противоречие.
        %
\end{proof}

\begin{problem}
        %
Условие конечности или счётности сигнатуры в этой теореме можно
ослабить. Как это сделать?
        %
\end{problem}

Вот пример применения теоремы~\ref{general-vaught}. Теория
алгебраически замкнутых полей характеристики~$0$ категорична в
любой несчётной мощности. (Это можно доказать, используя
базисы трансцендентности: такое поле имеет базис
трансцендентности над полем алгебраических чисел, мощность
которого равна мощности всего поля, а два поля с равномощными
базисами трансцендентности изоморфны). Следовательно, эта теория
полна.

Заметим, что это наблюдение согласовано со знаменитой (и
трудной!) теоремой Морли;%
\index{Теорема!Морли}%
\glossary{\Морли}
эта теорема утверждает, что теория с равенством,
категоричная в одной несчётной мощности, категорична и во всех
несчётных мощностях. (Подробно о теореме Морли можно прочесть,
например, в учебнике Кейслера\glossary{\Кейслер}
и Чэна\glossary{\Чэн}~\cite{keisler-chang}.)

\begin{theorem}
        \label{complete-decidable}%
Конечно аксиоматизируемая полная теория в конечной сигнатуре
разрешима.
        %
\end{theorem}

\begin{proof}
        %
Пусть дана произвольная формула $\varphi$. Будем перебирать все
выводы в исчислении предикатов и проверять, не обнаружилась ли
выводимость одной из формул~$\varphi$ или~$\lnot\varphi$ из
конъюнкции некоторых аксиом теории~$T$. Рано или поздно одна из
них окажется выводимой (поскольку теория полна), и тем самым мы
узнаем, какая из формул выводима в теории.
        %
\end{proof}

Это доказательство неконструктивно в том смысле, что не даёт
никакой оценки на время работы алгоритма. Отметим также, что не
обязательно требовать конечной аксиоматизируемости теории;
достаточно, чтобы она имела разрешимое или перечислимое
множество аксиом (см.~\cite{computable-functions}).

Проиллюстрируем все эти понятия на нескольких (в основном уже
обсуждавшихся нами) примерах.

\subsection*{Плотные линейно упорядоченные множества}

Рассмотрим сигнатуру, содержащую
отношения порядка и равенства. Рассмотрим \emph{теорию плотных
линейно упорядоченных множеств без первого и последнего
элемента},%
\index{Теория!плотных линейно упорядоченных множеств без первого и
последнего элемента|defin}
которая включает в себя следующие аксиомы:
        %
\begin{itemize}
        \item
аксиомы равенства (в том числе
сохранение порядка при замене элементов на равные);
        \item
$\forall x\,(x\le x)$
(рефлексивность порядка);
        \item
$\forall x\forall y\forall z\, ((x\le y)\land (y\le z)\to (x\le z))$
\\(транзитивность порядка);
        \item
$\forall x \forall y\,((x\le y)\land (y\le x)\to (x=y))$
\\(антисимметричность порядка);
           \item
$\forall x \forall y \,((x\le y)\lor (y\le x))$
(линейность порядка);
        \item
$\forall x \exists y\, (y >x)$ (нет максимального элемента;
$(y>x)$ можно считать сокращением для $\lnot (y\le x)$ или
для $(x\le y)\land\lnot (x=y)$\т при наличии остальных
аксиом это одно и то же);
        \item
аналогичная аксиома про отсутствие минимального элемента;
        \item
$\forall x\forall y\, ((x<y)\to \exists z\,((x<z)\land (z<y))$
(плотность).
        %
\end{itemize}

Рациональные числа образуют счётную модель этой
теории, а действительные\т несчётную.
Как мы уже упоминали, эта теория категорична в счётной мощности,
все её счётные нормальные модели изоморфны. Отсюда по
теореме~\ref{categorical-complete} получаем, что она полна.
Следовательно, в ней выводятся все истинные в $\mathbb{Q}$ (или
в любой другой модели, в частности, в $\mathbb{R}$) формулы её
сигнатуры (в самом деле, из формул $\varphi$ и $\lnot\varphi$
ровно одна истинна и ровно одна выводима, и выводимая формула
должна быть истинной). Наконец, по
теореме~\ref{complete-decidable} эта теория разрешима.

Другое доказательство тех же фактов даёт элиминация кванторов%
      \index{Элиминация кванторов}
(теорема~\ref{quantifier-elimination-dense},
с.~\pageref{quantifier-elimination-dense}). Как мы отмечали в
разделе~\ref{quantifier-elimination}, для каждой формулы
$\varphi$ нашей сигнатуры существует бескванторная
формула~$\varphi'$, эквивалентная $\varphi$ в любой нормальной
интерпретации теории плотных линейно упорядоченных множеств без
первого и последнего элементов. Поэтому эквивалентность
$\varphi\leftrightarrow\varphi'$ (с кванторами всеобщности)
является теоремой этой теории. Если формула $\varphi$ была
замкнутой, то формула $\varphi'$ будет тождественно истинной или
тождественно ложной. В первом случае в теории выводима
формула~$\varphi$, во втором случае\т её отрицание.
Следовательно, теория полна.

Сказанное можно интерпретировать и так: мы доказали конечную
аксиоматизируемость теории $\Th(\mathbb{Q},{=},{<})$,
предъявив список аксиом.

\begin{problem}
        %
Покажите, что эта теория не является категоричной в мощности
континуум.
        %
\end{problem}

Отсюда следует (по теореме Морли),%
\index{Теорема!Морли}%
\glossary{\Морли}
что теория плотных линейно
упорядоченных множеств без первого и последнего элемента не
будет категоричной ни в какой несчётной мощности.

\begin{problem}
        %
Не используя теоремы Морли, укажите примеры неизоморфных
плотных линейно упорядоченных множеств заданной несчётной
мощности.
        %
\end{problem}

\begin{problem}
        %
Покажите, что  теории $\Th([0,1],{=},{<})$ и
$\Th([0,+\infty),{=},{<})$ конечно аксиоматизируемы, полны и
разрешимы. Будут ли они категоричными в мощности континуум?
        %
\end{problem}

\begin{problem}
        %
Рассмотрим теорию плотных линейно упорядоченных множеств (не
добавляя аксиом про наименьший и наибольший элемент). Будет ли
она категорична в какой\д либо мощности? полна? разрешима?
        %
\end{problem}

\subsection*{Теория $\Th(\mathbb{Z},{=},{S},0)$}
        \index{Теория!$\Th(\mathbb{Z},{=},{S},0)$}%
В этом примере мы действуем в обратном порядке, начав с
конкретной интерпретации (целые числа с равенством, функцией
прибавления единицы и константой~$0$) и построив явную систему
аксиом. Для этого вспомним процедуру элиминации кванторов из
раздела~\ref{quantifier-elimination}
(теорема~\ref{quantifier-elimination-integers-successor}).
Какими свойствами должна обладать нормальная интерпретация
языка, чтобы преобразования, использованные при элиминации
кванторов, были эквивалентными? Помимо аксиом равенства, нам
нужно, чтобы функция~$S$ была биекцией и чтобы для любого $x$
все элементы
      $$
\dots,S^{-1}(S^{-1}(x)), S^{-1}(x), x,S(x),S(S(x)),\dots
      $$
были различны. Другими словами, элиминация кванторов даёт
формулу, эквивалентную исходной во всех моделях такой теории:
        %
\begin{itemize}
        \item
аксиомы равенства;
        \item
$\forall x \forall y\, ((S(x)=S(y))\to(x=y))$;
        \item
$\forall x \exists y\, (S(y)=x)$;
        \item
$\forall x \, \lnot (x=S(x))$;
        \item
$\forall x \, \lnot (x=S(S(x)))$;
        \item
$\forall x \, \lnot (x=S(S(S(x))))$ \итд
        %
\end{itemize}

Ограничиваясь замкнутыми формулами, мы (как и в предыдущем
примере) видим, что $\Th(\mathbb{Z},{=},{S},0)$ совпадает с
множеством всех формул, выводимых из перечисленных аксиом, так
что теория с этими аксиомами полна. Она разрешима (как любая
полная теория с разрешимым множеством аксиом; кроме того, явный
алгоритм даётся элиминацией кванторов).

В отличие от предыдущего примера, эта теория не является
категоричной в счётной мощности \т например,
$\mathbb{Z}+\mathbb{Z}$ является её моделью, не изоморфной
исходной.

\begin{problem}
        %
Опишите все модели этой теории.
        %
\end{problem}

\begin{problem}
        %
Покажите, что она категорична в любой несчётной мощности.
        %
\end{problem}

Покажем в заключение, что эта теория не является конечно
аксиоматизируемой. В самом деле, пусть имеется конечное
множество $F$ теорем этой теории, из которой следуют все
остальные теоремы. Каждая из теорем множества~$F$ выводима из
аксиом, и этот вывод использует конечное число аксиом. Это
означает, что все остальные аксиомы (не используемые в выводе
формул из~$F$) вообще лишние. А это не так: в конечной модели,
составленной из остатков по модулю~$N$, верны все аксиомы,
кроме $S^{N}(x)\ne x,\,S^{2N}(x)\ne x,\dots$, поэтому эти аксиомы
не следуют из остальных.

\begin{problem}
        %
Докажите, что использованное нами рассуждение носит общий характер:
если теория бесконечна, но конечно аксиоматизируема, то
некоторая её конечная часть равносильна всей теории (имеет те же
теоремы).
        %
\end{problem}

\subsection*{Теория $\Th(\mathbb{Z},{=},{<},{S},0)$}
\index{Теория!$\Th(\mathbb{Z},{=},{<},{S},0)$}%

Что изменится, если мы добавим к сигнатуре, помимо прибавления
единицы, ещё и отношение порядка? Как мы видели (см.\
доказательство теоремы \ref{successor-order-elimination} и
задачу после него, с.~\pageref{successor-order-elimination}),
элиминация кванторов по\д прежнему возможна. Для придания
законности нам нужны такие свойства интерпретации (которую мы
предполагаем нормальной): она представляет собой линейно
упорядоченное множество, в котором каждый элемент имеет
непосредственно следующий (совпадающий с значением функции~$S$)
и непосредственно предшествующий. В отличие от предыдущего
примера, нам достаточно конечного набора аксиом. Таким образом,
теория $\Th(\mathbb{Z},{=},{<},{S},0)$ конечно аксиоматизируема,
а также (как и в предыдущем примере) полна, разрешима, но не
категорична в счётной мощности.

Можно обойтись и без элиминации кванторов, рассуждая иначе.
Рассмотрим теорию линейно упорядоченных множеств со
следующим и предыдущим элементом и опишем все её модели. Именно,
мы покажем, что любая нормальная
модель~$M$ этой теории имеет вид
$\mathbb{Z}\times A$, где $A$\т произвольное линейно
упорядоченное множество (порядок на парах таков: сначала
сравниваются $A$\д компоненты, а в случае равенства\т
$\mathbb{Z}$\д компоненты.) В самом деле, будем говорить, что
элементы $x$ и $y$ лежат \лк в одной галактике\пк, если между
ними конечное число элементов. (Легко проверить, что это
действительно отношение эквивалентности, и наше множество
разбивается на галактики.)  Далее проверяем, что каждая
галактика изоморфна~$\mathbb{Z}$ (как упорядоченное множество)
и что на галактиках естественно определяется порядок.

Теперь с помощью игры Эренфойхта
(см.~раздел~\ref{elementary-equivalence},
теорема~\ref{z-equivalent-z-plus-z}) мы показываем, что все
нормальные модели этой теории элементарно эквивалентны. Отсюда
заключаем, что теория полна (как в доказательстве
теоремы~\ref{categorical-complete}, где мы по существу
использовали элементарную эквивалентность моделей, а не их
изоморфизм).

\begin{problem}
        %
Покажите, что теория $\Th(\mathbb{Z},{=},{<},S,0)$ не
категорична ни в какой несчётной мощности.
        %
\end{problem}

\begin{problem}
        %
Будет ли теория $\Th(\mathbb{Z},{=},{<})$ конечно аксиоматизируемой?
разрешимой? категоричной?
        %
\end{problem}

\begin{problem}
        %
Будет ли теория $\Th(\mathbb{N},{=},{<})$ конечно аксиоматизируемой?
разрешимой? категоричной?
        %
\end{problem}

\subsection*{Теория $\Th(\mathbb{Q},{=},{<},{+},0,1)$}
\index{Теория!$\Th(\mathbb{Q},{=},{<},{+},0,1)$}%
Эту теорию мы рассматривали в
разделе~\ref{quantifier-elimination},
с.~\pageref{q-plus-order-elimination}. Мы ограничимся двумя
константами~$0$ и~$1$, поскольку любую атомарную формулу можно
привести к общему знаменателю и получить целые константы,
которые можно выразить через~$0$ и~$1$.

Мы хотим указать явно набор аксиом этой теории, то есть
множество формул, из которых выводятся все теоремы этой теории и
только они. Как и в предыдущих примерах, это можно сделать,
проанализировав процесс элиминации кванторов и выявив все
использованные при этом свойства интерпретации. (Все
рассматриваемые нами интерпретации предполагаются нормальными, а
аксиомы равенства изначально включаются в строимую нами теорию.)

Прежде всего, нам важно, что по сложению мы имеем абелеву группу
(и $0$ является её нулём). Это позволяет в равенствах переносить
члены с одной стороны в другую. Для операций с неравенствами нам
надо знать, что порядок является линейным и что он согласован со
сложением (то есть что из $x\hm<y$ следует~${x+z}\hm<{y+z}$).
Кроме того, мы умножали равенства и неравенства на рациональные
числа. Чтобы это было законно, мы должны знать, что группа
является делимой: для всякого $a$ уравнения ${x+x}\hm=a,\,
{x+x+x}\hm=a,\,{x+x+x+x}\hm=a,\,\dots$ имеют решения. (В
упорядоченной группе такое решение, как легко показать,
единственно.) Наконец, нам надо знать, что~$1>0$.

Кроме этих аксиом (которых счётное число) мы при элиминации
ничего не использовали, так что для любой формулы $\varphi$ есть
бескванторная формула $\varphi'$, которая эквивалентна $\varphi$
в любой делимой упорядоченной группе. Поэтому любая замкнутая
формула, истинная в стандартной интерпретации (в~$\mathbb{Q}$),
истинна в любой делимой упорядоченной группе, и мы получили
счётную систему аксиом для теории
$\Th(\mathbb{Q},{=},{<},{+},0,1)$.

\begin{problem}
        %
Покажите, что эта теория не является конечно аксиоматизируемой.
(Указание: делимость любого элемента группы
на простое число~$p$ не вытекает из делимости
на все меньшие простые числа\т рассмотрим рациональные числа,
знаменатель которых взаимно прост с $p$.)
        %
\end{problem}

\begin{problem}
        %
Покажите, что эта теория разрешима.
        %
\end{problem}

\begin{problem}
        %
Покажите, что эта теория не является категоричной.
        %
\end{problem}

\begin{problem}
        %
Покажите, что теория $\Th(\mathbb{Q},{=},{+},0)$ не является
категоричной в счётной мощности, но категорична в любой несчётной
мощности. (Указание: её модели\т векторные пространства над
полем рациональных чисел.)
        %
\end{problem}

\subsection*{Арифметика Пресбургера}
\index{Арифметика!Пресбургера}%
\glossary{\Пресбургер}%
В разделе~\ref{presburger} мы занимались элиминацией кванторов в
теории $(\mathbb{Z},{=},{<},{+},{0},{1})$, которая потребовала
добавления бесконечного числа дополнительных предикатов
(сравнимость по модулю $N$ для всех целых $N>1$).

Проанализировав это рассуждение, можно извлечь из него явную
аксиоматизацию для теории $(\mathbb{Z},{=},{<},{+},{0},{1})$
(без дополнительных предикатов). Какие свойства порядка и
сложения на целых числах мы используем? Нам важно, что целые
числа образуют абелеву группу, что порядок согласован со
сложением и что~${x+1}$ есть непосредственно следующий за~$x$
элемент (достаточно, впрочем, сказать, что $1$ непосредственно
следует за~$0$).
В любой группе можно рассмотреть подгруппу делящихся
на~$N$ элементов (для любой целой константы $N>1$) и сравнивать
элементы по модулю этой подгруппы. Но этого мало: нам нужно ещё
иметь возможность делить на $N$ с остатком. Это гарантируется
такой аксиомой (при каждом~$N$\т своя аксиома): для любого
элемента $x$ ровно один из $N$ элементов $x, {x-1},
{x-2},\dots,{x-N+1}$ делится на $N$.

Можно проверить, что все шаги элиминации кванторов сохраняют
равносильность в такой ситуации. Проверим, например, что
сравнения можно умножать на целое положительное число. Почему,
скажем, ${a\equiv b \pmod 3}$ равносильно
${a+a}\hm\equiv{b+b\pmod 6}$? По определению первое означает,
что $a-b=u+u+u$ для некоторого~$u$, а второе\т что
${(a-b)}\hm+{(a-b)}\hm=v\hm+v\hm+v\hm+v\hm+v\hm+v$ для
некоторого~$v$, и достаточно сослаться на то, что в
упорядоченной группе из ${x+x}\hm=0$ следует $x\hm=0$ (поскольку
из $x\hm>0$ следует $x\hm+x\hm>x\hm>0$). Наиболее сложный шаг\т
доказательство представительности набора. Здесь надо рассмотреть
все случаи расположения произвольного $y$ относительно правых
частей. В каждом из случаев мы заменяли $y$ на первый элемент из
окрестности правых частей, встречающийся при движении шагами $D$
(как описано на с.~\pageref{presburger-set}). Эту процедуру
можно понимать как деление расстояния (до ближайшей правой
части) на~$D$ с остатком; возможность этого гарантируется нашими
аксиомами.

\begin{problem}
        %
Покажите, что теория
$(\mathbb{Z},{=},{<},{+},{0},{1})$ разрешима.
        %
\end{problem}

\begin{problem}
        %
Покажите, что эта теория не категорична в счётной мощности.
(Указание: Среди её моделей есть модели вида
$\mathbb{Z}\times A$, где $A$\т делимая упорядоченная группа.\footnote{%
        В первом издании книги ошибочно утверждалось, что все модели
        этой теории имеют такой вид. Авторы благодарны Евгению Петровичу
        Бондаренко, который обнаружил эту ошибку и указал контрпример.}%
)
        %
\end{problem}

\begin{problem}
        %
Покажите, что эта теория не конечно аксиоматизируема. (Указание:
рассмотрите интерпретацию $\mathbb{Z}\times A$, когда в $A$
возможно деление на простые числа, меньшие некоторого~$p$, но не
на~$p$.)
        %
\end{problem}

\subsection*{Алгебраически замкнутые поля характеристики~$0$}
 \index{Теория!алгебраически замкнутых полей характеристики~$0$}%
 \index{Поле!алгебраически замкнутое}%
 \index{Алгебраически замкнутое поле}%
 \index{Характеристика поля}%

Теорема~\ref{c-quantifier-elimination} устанавливает
возможность элиминации кванторов в поле комплексных чисел. Как
 \index{Комплексные числа}%
мы уже отмечали
(теорема~\ref{algebraically-closed-fields-completeness} на
с.~\pageref{algebraically-closed-fields-completeness}), это
преобразование даёт формулу, равносильную во всех алгебраически
замкнутых полях характеристики~$0$. Отсюда следует, как обычно, что
теория алгебраически замкнутых полей характеристики~$0$ полна.
(Её аксиомы таковы: аксиомы равенства, аксиомы
поля, счётный набор утверждений о
том, что любой многочлен степени~$n\hm>0$ с ненулевым старшим
коэффициентом имеет корень, а также счётный набор утверждений
вида $1\hm+1\hm+1+\ldots+1\ne 0$.) Эта теория совпадает с
элементарной теорией поля комплексных чисел.

Другое доказательство полноты этой теории можно получить с
помощью критерия Лося\ч Воота (теорема~\ref{general-vaught},
с.~\pageref{general-vaught})
        \glossary{\Лось}%
        \glossary{\Воот}%
и утверждения
следующей задачи.

\begin{problem}
        %
Покажите, что теория алгебраически замкнутых полей
характеристики~$0$ не категорична в счётной мощности, но
        \index{Категоричная теория}%
        \index{Теория!категоричная}%
категорична в любой несчётной мощности. (Указание: алгебраически
замкнутое поле имеет базис трансцендентности над~$\mathbb{Q}$; если
        \index{Базис трансцендентности}%
поле несчётно, то базис равномощен полю.)
        %
\end{problem}

\begin{problem}
        %
Покажите, что теория алгебраически замкнутых полей данной
характеристики~$p$ полна.
        %
\end{problem}

\begin{problem}
        %
Покажите, что если некоторая формула в сигнатуре~$({=},{+},{\times})$
истинна в алгебраически замкнутых полях
характеристики~$0$, то она истинна и во всех алгебраически
замкнутых полях достаточно большой конечной характеристики.
        %
\end{problem}

\subsection*{Вещественно замкнутые поля}
        \label{real-closed-fields}%

Более сложно сформулировать, какие свойства поля вещественных чисел
реально используются при доказательстве теоремы Тарского\ч
Зайденберга (раздел~\ref{seidenberg-tarski}). Дело в том, что мы
ссылались на разные факты из анализа (понятие производной,
теорема Ролля, асимптотика многочленов). Однако на самом деле
достаточно некоторых алгебраических свойств поля\т как говорят,
поле должно быть \лк вещественно замкнутым\пк.

Поле называется \emph{упорядоченным},%
\index{Упорядоченное поле|defin}%
\index{Поле!упорядоченное|defin}
если на нём задан линейный
порядок, причём он согласован со сложением ($x<y$ влечёт
$x+z<y+z$) и умножением ($x,y>0$ влечёт $xy>0$).

\begin{problem}
        %
Покажите, что любое упорядоченное поле имеет характеристику~$0$
и что в нём сумма квадратов не может
равняться $-1$.
        %
\end{problem}

Упорядоченное поле называется \emph{вещественно замкнутым},%
\index{Вещественно замкнутое поле|defin}%
\index{Поле!вещественно замкнутое|defin}
если любой многочлен, имеющий на концах отрезка разные знаки, имеет
корень на этом отрезке. (Отметим в скобках, что существует
несколько эквивалентных определений вещественно замкнутого поля,
см.~учебник ван дер Вардена~\cite{vanderwaerden};
      \glossary{\вандерВарден}%
мы выбрали наиболее удобное для наших~целей.)

Мы не можем записать определение вещественной замкнутости в виде
формулы, поскольку степень многочлена может быть любой. Но можно
написать много формул\т по одной для каждой степени многочлена.
Например, для многочленов степени $2$ получится формула
\begin{multline*}
(u<v)\ \land\ (a\ne 0)\  \land\ ((au^2+bu+c)(av^2+bv+c)<0) \to\\
\to\exists x\, ((u<x)\land (x<v)\land\ (ax^2+bx+c=0)).
\end{multline*}

Теория, состоящая из аксиом упорядоченного поля (в том числе
аксиом равенства) и этих дополнительных аксиом, называется
\emph{теорией вещественно замкнутых полей}.%
    \index{Теория!вещественно замкнутых полей|defin}%
    \index{Поле!вещественно замкнутое|defin}
Покажем, что в любом вещественно замкнутом поле
выполнены основные факты о многочленах и их производных. Прежде
всего заметим, что в алгебре естественно определять производную
многочлена не как предел, а чисто формально: $(x^n)'=nx^{n-1}$
(для любого положительного целого~$n$), далее по линейности.
Степень производной многочлена на единицу меньше степени самого
многочлена. Выполнены основные правила дифференцирования
(линейность, правило дифференцирования произведения,
формула~Тейлора).%
    \index{Формула Тейлора}%
    \glossary{\Тейлор}

Теперь отметим некоторые свойства многочленов, связанные с
порядком. Пусть многочлен в какой\д то точке равен нулю, а
производная его в этой точке больше нуля. Тогда в некоторой
окрестности этой точки он положителен справа и отрицателен
слева. (В самом деле, можно применить формулу Тейлора и
оценки, показывающие, что вблизи нуля знак определяется
линейной частью.)

Справедливо и свойство сохранения знака: если в какой\д то точке
многочлен положителен, то и в достаточно близких точках он
положителен. Слова \лк достаточно близких\пк\ понимаются в
обычном смысле ($\exists \varepsilon > 0$), только теперь
$\varepsilon$\т не действительное число, а элемент поля.

Аналогичным образом можно сформулировать и доказать такое
утверждение: при всех достаточно больших значениях аргумента
знак многочлена определяется его старшим коэффициентом (обычные
оценки вполне годятся).

Более сложно доказывается, что что если производная многочлена
$P(x)$ положительна на интервале $(a,b)$, то он неубывает на
отрезке $[a,b]$. Пусть это не так. Добавив к многочлену
константу, можно считать, что для каких\д то точек $c,d$ этого
отрезка имеют место неравенства $c<d$, $P(c)\hm>0$,
$P(d)\hm<0$ и потому $P$ имеет корень на интервале~$(c,d)$
(вещественная замкнутость).

Число корней многочлена (в любом поле) конечно, поэтому среди
корней многочлена $P(x)$ на интервале $(c,d)$ есть наименьший
корень $\alpha$. Слева от $\alpha$ многочлен $P$ должен быть
положительным (поскольку корень первый), но одновременно вблизи
$\alpha$ и отрицательным (так как $P'(\alpha)\hm>0$ по
предположению).

Теперь легко понять, что многочлен с положительной производной
на~$(a,b)$ строго возрастает на~$[a,b]$: если бы в двух точках он
принимал одинаковые значения, то между ними он был бы константой
(чего не может быть для непостоянного многочлена).

Следствие: многочлен, производная которого не имеет корней на
$(a,b)$, либо строго возрастает, либо строго убывает на~$[a,b]$.
В самом деле, свойство вещественной замкнутости можно применить
и к производной, следовательно она либо всюду положительна, либо
всюду отрицательна. В частности, справедлива теорема Ролля%
\index{Теорема!Ролля}
\glossary{\Ролль}
(многочлен с равными значениями на концах отрезка имеет нуль
производной).

После такой тренировки наши рассуждения про знаки многочленов
диаграммы легко провести для произвольного поля. Легко понять
также, что деление с остатком (точнее, операция
модифицированного остатка) имеет смысл для любого поля, и что
целые числа (которые были коэффициентами многочленов) содержатся
в любом~поле.

Итак, элиминация кванторов даёт формулу, равносильную исходной в
любом вещественно замкнутом поле. Отсюда, как обычно, следует,
что теория вещественно замкнутых полей совпадает с элементарной
теорией упорядоченного поля вещественных чисел и потому полна, а
также разрешима, и что все вещественно замкнутые упорядоченные
поля элементарно эквивалентны.

% как бы доказать, что она не конечно аксиоматизируема?
